\documentclass[answers,12pt,addpoints]{exam} 
\usepackage{import}

\import{C:/Users/prana/OneDrive/Desktop/MathNotes}{style.tex}

% Header
\newcommand{\name}{Pranav Tikkawar}
\newcommand{\course}{Spatial Statistics}
\newcommand{\assignment}{Homework 4.3}
\author{\name}
\title{\course \ - \assignment}

\begin{document}
\maketitle

\newpage
\begin{questions}
\question 5.7

Derive 5.18: 
$$ \frac{\partial}{\partial \theta_j} \log | \Sigma (\theta) | = \trace\{ \Sigma^{-1} \Sigma_j(\theta) \} $$
\begin{solution}
    Consider an $m$ dimensional Gaussian random vector $Y \sim N(\mu(\beta), \Sigma(\theta))$.
    
    Using the Jacobi identity $\partial (\det A) = \det A \cdot \trace (A^{-1} \partial A)$, we have
    \begin{align*}
        \frac{\partial}{\partial \theta_j} \log | \Sigma (\theta) | &= \frac{1}{|\Sigma(\theta)|} \cdot \frac{\partial}{\partial \theta_j} |\Sigma(\theta)| \\
        &= \frac{1}{|\Sigma(\theta)|} \cdot |\Sigma(\theta)| \cdot \trace\{ \Sigma^{-1} \Sigma_j(\theta) \} \\
        &= \trace\{ \Sigma^{-1} \Sigma_j(\theta) \}
    \end{align*}

    
\end{solution}

\question 6.5 a

\begin{solution}
    For a log-log plot, the slope of the line is given by the exponent of the power law. For all $m =100,200,400$ we see that the slope is approximately $.5$ and when regraphing on a regular plot, we see that the line is approximately $y = \sqrt{x}$, which confirms our observation from the log-log plot.
\end{solution}

\question 6.5 b

\begin{solution}
    Let us consider the Taylor expansion of the function $K_\theta(x) = \theta_1 e^{-r / \theta_2}(1+r/\theta_2)$ around $r=0$:
\begin{align*}
    K_\theta(x) &= \theta_1 e^{-r / \theta_2}(1+r/\theta_2) \\
    &= \theta_1 (1 - r/\theta_2 + (r/\theta_2)^2/2 - o((r/\theta_2)^2))(1 + r/\theta_2) \\
    &= \theta_1 (1 + r/\theta_2 - r/\theta_2 + (r/\theta_2)^2/2 - o((r/\theta_2)^2)) \\
    &= \theta_1 (1 + (r/\theta_2)^2/2 - o((r/\theta_2)^2))
\end{align*}
Notice that our principle irregular term is of order $r^2$, with the coefficient $\theta_1 / (2 \theta_2^2)$. 

Going back to the plot of $\theta_1$ vs $\theta_2$, we see that the line of best fit on the log log scale has a slope of approximately $\frac{1}{2}$ and that meand the model that is fit is 
\begin{align*}
    \log \theta_2 &= \frac{1}{2} \log \theta_1 + b \\
    \theta_2 &= B \cdot \theta_1^{1/2} \\
    \frac{\theta_2}{\theta_1^{1/2}} &= B
\end{align*}
This finding is consistent with the fact that the coefficient of the principle irregular term is $\theta_1 / (2 \theta_2^2)$, which is constant when $\theta_2 / \theta_1^{1/2}$ is constant.

Thus our reparameterization is given by $\phi_1 = \theta_2 / \theta_1^{1/2}$ and $\phi_2 = \theta_2 \theta_1^{1/2}$.

Note that we could have easiley also chosen $\phi_1 = \theta_1 / \theta_2^2$ and $\phi_2 = \theta_1 \theta_2^2$, which would have also worked.
\end{solution}

\end{questions}

\end{document}